\begin{cabstract}
对冲基金交易是金融市场中高度系统化、自动化和博弈化的决策活动，其核心并非单一价格预测，而是贯穿资产选择、因子分析、时机识别、交易执行与最终交易判断的连续决策过程。在多空博弈、高杠杆约束、高维非平稳时序以及严格执行成本共同作用下，传统金融预测模型与单一对抗学习范式容易面临策略探索不足、模式崩溃、噪声过拟合和执行偏差放大等问题。围绕真实交易流程构建兼具鲁棒性、泛化性与可解释性的对抗学习方法，对于提升对冲基金在复杂市场环境下的风险收益表现具有重要理论意义与应用价值。

在真实交易场景中，对冲基金交易决策通常不是一次性给出买卖信号，而是沿资产配置依据、因子解释依据、时序结构依据、执行路径约束和最终策略整合逐层展开。具体而言，资产配置环节需要在多资产池中识别相对配置价值并确定资金权重；因子建模环节需要从高维金融变量中筛选具有稳定解释能力的信息；时序结构识别环节需要在非平稳行情中判断趋势形态和交易时机；交易执行环节需要在订单落地过程中控制滑点、市场冲击和执行偏差；策略整合环节则需要将方向判断、结构信号和执行约束转化为买入、卖出或持有等交易动作。现有方法在这些环节中仍存在多资产配置关系不稳定、有效因子易受噪声干扰、时序结构识别受局部伪结构影响、执行层风险难以内生控制、最终交易判断约束不足等问题。这些问题在真实交易中相互传导，使对抗学习方法的应用仍面临较高难度。

为解决上述问题，本文不再将对冲基金交易算法视为单一模型优化问题，而是将其理解为由投资组合管理、高维金融因子建模、时序结构识别、交易执行优化和策略整合共同构成的完整交易决策流程。围绕这一流程，本文提出多类面向对冲基金交易关键环节的多智能体对抗学习方法，使各章研究内容分别服务于交易决策链条中的具体问题。本文主要贡献如下。

第一，在资产配置与组合管理环节，提出面向投资组合管理的群组交叉对抗（Group Cross Adversarial, GCA）学习方法。针对单一对手条件下策略探索空间受限、多资产相对关系刻画不稳定以及组合权重生成易受局部信号扰动影响的问题，该方法将智能体划分为若干群组并引入群组间交叉对抗，形成有序的群组交叉博弈结构，用于刻画多资产之间的资金配置关系，并抑制单一博弈路径下的过拟合。本文以多品种期货与股票组合交易任务作为验证场景开展实验。实验显示，该方法在夏普比率与最大回撤控制方面优于传统基准模型，说明其能够服务于多资产资金配置和组合风险控制。

第二，在交易依据解释与高维金融因子建模环节，提出面向高维金融因子建模的赛马场竞争方法。针对高维因子空间中有效信息容易被噪声掩盖、模型表征容易受到非对称竞争压制并出现过早收敛的问题，该方法通过引入多异构模型在统一评分规则下的动态竞争，构建成熟模型向新生模型传递稳定表征、新生模型以差异化反馈反向约束成熟模型的双向知识迁移逻辑，从而筛选并强化具有稳定解释能力的金融因子，延迟单一模型的过早收敛。本文在复杂耦合因子模拟环境与真实高频数据上开展验证实验。实验显示，赛马场竞争方法能够增强高维因子表示的鲁棒性，并缓解因子提取过程中的信息丢失问题。

第三，在交易时机判断与时序结构识别环节，提出面向时序结构识别的逆向师生蒸馏学习方法。针对主观波浪形态难以量化、传统“强教弱”范式在非平稳噪声下容易陷入局部最优以及结构信号交易意义难以判别的问题，该方法结合四元分形事件对艾略特波浪理论进行量化重构，并设计从弱模型向强模型传递失败经验的逆向蒸馏路径，促使最优模型向组内分布靠拢，以抑制噪声过拟合。本文在具有复杂趋势形态的行情场景下开展验证实验。实验结果显示，该方法在结构识别效果上优于主流基线模型，能够滤除部分高频白噪声并保留趋势结构信息，为交易时机识别提供支撑。

第四，在执行路径优化与交易执行控制环节，提出面向交易执行优化的孪生对抗自体再生方法。针对交易信号进入真实成交环节后可能面临的成交量预测偏差、执行路径失衡、滑点风险与对抗崩溃问题，该方法通过“复制—配对—分化”的自体修复逻辑，在组内对抗发生崩溃时自动触发跨组修复程序，将市场冲击风险内化为可优化的数学目标。本文以成交量加权平均价格执行策略作为验证对象，在真实交易数据上开展实验。实验结果显示，该方法在执行滑点控制方面具有一定优势，可为细粒度交易执行优化提供崩溃修复思路。

第五，在最终交易判断与策略整合环节，提出面向策略整合的偏序协同对抗（Partial-order Collaborative Adversarial, PCA）学习方法。针对金融市场层级博弈关系、模式崩溃以及最终交易判断中多源约束难以协调的问题，该方法采用非对称配置模拟不同资金体量参与者之间的博弈关系，并将涨、跌、平分类结果解释为买入、卖出或持有等交易动作。在表述口径上，本章承接资产配置、因子依据、结构信号与执行约束，但不将其写成对前述实验结果的直接拼接。本文选取具备非对称博弈特征的不同资产类别开展对比验证。实验结果显示，该方法具有较好的跨资产泛化能力，并改善了最终判断环节的稳定性。

围绕上述研究内容，本文沿对冲基金交易决策链条提出多类多智能体对抗学习方法：第2章以 GCA 支撑多资产资金配置，第3章构建有效金融因子的动态识别机制，第4章面向趋势结构与交易时机进行时序结构识别，第5章围绕订单落地过程优化交易执行，第6章以 PCA 将涨跌分类组织为最终交易判断与策略整合问题。上述研究从资产配置、因子筛选、结构识别、执行优化到最终策略整合逐步展开，能够有效应对高维非平稳时序数据、复杂博弈结构及极端行情下的建模挑战，为对冲基金交易提供更加稳健、高效且可解释的决策支持。
\end{cabstract}

\ckeywords{对冲基金交易,多智能体对抗学习,群组交叉对抗,偏序协同对抗,策略整合}
